\subsubsection{杜教筛 (Du's Sieve)}
杜教筛，用于在亚线性时间内求解数论函数（如 $\mu, \phi$）的前缀和。
设 $S(n) = \sum_{i=1}^n f(i)$，我们构造 $g$ 使得 $\sum (f*g)$ 和 $\sum g$ 易于计算。
其核心公式为：
$$ g(1)S(n) = \sum_{i=1}^n (f*g)(i) - \sum_{i=2}^n g(i)S(\lfloor n/i \rfloor) $$
\begin{lstlisting}
ll n,m;
ll pr[2000010];
bool vis[2000010];
ll mu[2000010];
ll ph[2000010];
void gt(){
    static bool done=0;if(done)return;done=1;cnt=0;
	mu[1]=ph[1]=1;
	f(i,2,2000000){
		if(!vis[i])pr[++cnt]=i,mu[i]=-1,ph[i]=i-1;
		for(int j=1;j<=cnt&&i*pr[j]<=2000000;j++){
			vis[i*pr[j]]=1;if(i%pr[j]==0){
				ph[i*pr[j]]=ph[i]*pr[j];break;
			}mu[i*pr[j]]=-mu[i],ph[i*pr[j]]=ph[i]*ph[pr[j]];
		}
	}f(i,2,2000000)mu[i]+=mu[i-1],ph[i]+=ph[i-1];
}
unordered_map<ll,ll>smu;
unordered_map<ll,ll>sph;
ll Mu(ll n){
	if(n<=2000000)return mu[n];
	if(auto it=smu.find(n);it!=smu.end())return it->second;
	ll res=1;for(ll l=2,r;l<=n;l=r+1){
		r=n/(n/l);res-=(r-l+1)*Mu(n/l);
	}return smu[n]=res;
}ll Ph(ll n){
	if(n<=2000000)return ph[n];
	if(auto it=sph.find(n);it!=sph.end())return it->second;
	ll res=n*(n+1)/2;for(ll l=2,r;l<=n;l=r+1){
		r=n/(n/l);res-=(r-l+1)*Ph(n/l);
	}return sph[n]=res;
}
int main(){
	gt();wt{
		n=d;printf("%lld %lld\n",Ph(n),Mu(n));
	}
	return 0;
}
\end{lstlisting}

\paragraph{多测与注意} gt 一次，前缀表和 memo 可跨组保留。只修正“缓存值为 0 时重复计算”；筛上限按题目调整。
